% Source: Solutions/Ansh-PS1-Q1-5.pdf, pp. 1-2 (typed). Restated problem statement omitted.
% REVIEW: the remark "including A = emptyset ... so no map g: emptyset -> P(emptyset) exists at
%         all" is wrong -- there is exactly one map (the empty function); it is simply not
%         surjective, since P(emptyset) = {emptyset} is non-empty. Kept as written.
\submittedby{Ansh}

\begin{proof}
We establish both claims rigorously using Cantor's diagonalization argument, and we also show
how the non-existence of an injection follows directly from the non-existence of a surjection.

\paragraph{Part I: Non-existence of a surjective map $g\colon A\to\Pow(A)$.}
Suppose, for the sake of contradiction, that there exists a surjective function
$g\colon A\to\Pow(A)$. Define the diagonal subset $B\subseteq A$ by
\[ B:=\{x\in A\mid x\notin g(x)\}. \]
Note that this construction is well-defined for every set $A$ (including $A=\emptyset$, in
which case $\Pow(\emptyset)=\{\emptyset\}\neq\emptyset$, so no map
$g\colon\emptyset\to\Pow(\emptyset)$ exists at all).

Since $B\subseteq A$, we have $B\in\Pow(A)$. Because $g$ is assumed to be surjective, there
exists at least one element $x_0\in A$ such that
\[ g(x_0)=B. \]
By the law of excluded middle, either $x_0\in B$ or $x_0\notin B$. We examine both cases:
\begin{enumerate}[label=(\alph*)]
  \item \emph{Case 1 ($x_0\in B$):} By the defining property of $B$, an element $x\in A$
        belongs to $B$ if and only if $x\notin g(x)$. Consequently,
        $x_0\in B\implies x_0\notin g(x_0)$. However, since $g(x_0)=B$, this implies
        $x_0\notin B$, contradicting our assumption that $x_0\in B$.
  \item \emph{Case 2 ($x_0\notin B$):} Since $g(x_0)=B$, the assumption $x_0\notin B$ means
        $x_0\notin g(x_0)$. By the definition of $B$, any element of $A$ satisfying
        $x_0\notin g(x_0)$ must belong to $B$; hence $x_0\in B$, contradicting our assumption
        that $x_0\notin B$.
\end{enumerate}
In both cases we arrive at the logical contradiction $x_0\in B\iff x_0\notin B$. Therefore,
our initial assumption is false, and no surjective map $g\colon A\to\Pow(A)$ exists.

\paragraph{Part II: Non-existence of an injective map $f\colon\Pow(A)\to A$.}
We provide two independent proofs of this result.

\textbf{Method 1 (Direct Diagonal Proof).} Suppose, for the sake of contradiction, that there
exists an injective map $f\colon\Pow(A)\to A$. Consider the subset $C\subseteq A$ defined by
\[ C:=\{f(S)\in A\mid S\in\Pow(A)\text{ and }f(S)\notin S\}. \]
Since $C\subseteq A$, we have $C\in\Pow(A)$. Let $c:=f(C)\in A$. We ask whether $c\in C$ or
$c\notin C$:
\begin{enumerate}[label=(\alph*)]
  \item \emph{Case 1 ($c\in C$):} By definition of $C$, there exists some subset
        $S\in\Pow(A)$ such that $c=f(S)$ and $f(S)\notin S$. Since $f(C)=c=f(S)$ and $f$ is
        injective, it follows that $C=S$. Substituting $S=C$ into $f(S)\notin S$ yields
        $c=f(C)\notin C$, a contradiction.
  \item \emph{Case 2 ($c\notin C$):} Since $c=f(C)$, the statement $c\notin C$ is equivalent to
        $f(C)\notin C$. Taking $S=C\in\Pow(A)$ in the defining condition for $C$, we see that
        $f(C)$ satisfies the membership criterion for $C$, which implies $c=f(C)\in C$, a
        contradiction.
\end{enumerate}
Thus, in both cases we obtain a contradiction. Hence, no injective map $f\colon\Pow(A)\to A$
can exist.

\textbf{Method 2 (Reduction to Part I).} Suppose $f\colon\Pow(A)\to A$ is injective. Since
$\Pow(A)$ is always non-empty (as $\emptyset\in\Pow(A)$), its image $f(\Pow(A))\subseteq A$ is
non-empty, which implies $A\neq\emptyset$. Define a map $g\colon A\to\Pow(A)$ by
\[ g(a):=\begin{cases}
  S & \text{if } a=f(S) \text{ for some } S\in\Pow(A),\\
  \emptyset & \text{if } a\notin f(\Pow(A)).
\end{cases} \]
Because $f$ is injective, for each $a\in f(\Pow(A))$ the preimage $S\in\Pow(A)$ with $f(S)=a$
is unique; thus $g$ is a well-defined function. Moreover, for every $S\in\Pow(A)$, we have
$g(f(S))=S$, showing that $g\colon A\to\Pow(A)$ is surjective. This directly contradicts
Part I. Therefore, no injective map $f\colon\Pow(A)\to A$ exists.
\end{proof}
